% =============================================================
% ECCV Introduction — SOMO
% =============================================================
% Prerequisites: \newcommand{\modelname}{\texttt{SOMO}\xspace}
% Citations needed (placeholders):
%   - Object-centric: li2023gendexgrasp, chen2024multihandgrasp,
%     liu2023contactgen, jiang2021graspTTA
%   - Interaction-centric: chen2024dro (D(R,O)), chen2024tro (T(R,O))
%   - Hand-specific: xu2023unidexgrasp, wan2023unidexgrasp++
%   - Diffusion models: ho2020denoising, song2020denoising
%   - Classifier-free guidance: ho2022classifierfree
%   - TripoSG: triposg
% =============================================================

\section{Introduction}
\label{sec:intro}

Dexterous grasping is a fundamental capability for achieving precise, human-level manipulation.
Yet, generating grasps for \emph{any} robot hand---from a two-fingered gripper to a twenty-one-link anthropomorphic hand---remains an open problem.
While individual hands can be addressed by dedicated models, the growing diversity of robotic end-effectors demands a \emph{cross-embodiment} solution: a single generative model that synthesizes high-quality grasps across morphologically distinct hands.

Existing learning-based approaches to cross-embodiment dexterous grasping can be broadly categorized into two paradigms.
%
\textbf{Object-centric methods}~\cite{li2023gendexgrasp,chen2024multihandgrasp,liu2023contactgen} describe grasps in terms of object-surface features such as contact maps or affordance fields.
By grounding the representation in object geometry, these approaches naturally promote transferability across different hands and object shapes.
However, they are often brittle under partial observations and typically require an additional optimization stage---such as solving fingertip inverse kinematics or fitting predicted contact maps under penetration-free and joint-limit constraints---to convert abstract predictions into feasible joint configurations, making the pipeline computationally expensive and difficult to scale.
%
\textbf{Interaction-centric methods} move beyond pure object-surface descriptions by jointly modeling the spatial relationship between the hand and the object.
D(R,\,O) Grasp~\cite{chen2024dro} pioneered this direction, constructing a distance matrix between points sampled on the hand and the object surface to embed both entities in a shared interaction space.
This formulation retains embodiment awareness while capturing transferable object-level features, facilitating generalization across both hands and objects.
However, the point-to-point feature maps in D(R,\,O) demand substantial GPU memory, making the representation resource-intensive and difficult to scale.
Moreover, it relies on a suitable initial configuration, rendering performance sensitive to initialization quality.
T(R,\,O) Grasp~\cite{chen2024tro} addresses these limitations by replacing the dense distance matrix with a heterogeneous graph of SE(3) transforms between robot links and object patches, coupled with a graph diffusion model that generates grasps from pure noise without requiring initialization.
While T(R,\,O) demonstrates strong results on individual hands, it trains a separate model per embodiment, inheriting the fundamental scalability bottleneck: every new hand requires its own paired object-grasp dataset and its own training run, with no knowledge transfer between hands.

We argue that a truly scalable framework for cross-embodiment dexterous grasping must satisfy two key properties.
First, it should employ a \emph{unified representation} that naturally accommodates hands with heterogeneous kinematic structures---varying numbers of joints, link topologies, and coupling constraints---within a single generative model, without any per-hand architectural components.
Second, its performance should \emph{scale positively} with the number of training embodiments: adding more hands to the training pool should improve---or at the very least preserve---grasp quality for every individual hand, rather than degrading it through interference.

We present \modelname, a Morphology-Prior Diffusion model that realizes both properties through three synergistic designs---each enabling the next.
%
\textbf{(i)~Grasping as 3D assembly.}
Rather than predicting hand-specific joint angles, \modelname decomposes each robot hand into its constituent rigid links and predicts a per-link SE(3) pose.
A valid grasp becomes an arrangement of rigid pieces around the target object---akin to assembling a 3D puzzle---with kinematic feasibility enforced only \emph{after} generation through inverse kinematics.
Because the generative model operates on variable-length sets of SE(3) tokens rather than fixed-dimensional joint vectors, a single transformer-based denoiser can natively process embodiments ranging from 5-link grippers to 21-link anthropomorphic hands without any architecture modification.
This kinematics-agnostic representation is what makes the next two designs possible.
%
\textbf{(ii)~Shared 3D shape-aware encoding.}
A frozen pretrained 3D VAE encodes both object point clouds and individual link meshes into a common geometric latent space using the same encoder weights.
Object tokens capture local surface geometry through spatially distributed patch embeddings; morphology-aware link embeddings encode each link's shape, size, and contact surface.
This shared space lets the denoiser reason about how a link's geometry complements the local object surface, and maps novel link geometries into the same embedding space as known hands.
%
\textbf{(iii)~Classifier-free training with a morphology prior.}
We extend the classifier-free guidance paradigm~\cite{ho2022classifierfree} beyond its original label-dropping formulation: the unconditional path is trained on a distinct \emph{morphology-prior distribution} that captures physically valid hand poses through forward-kinematic self-exploration, without any object-grasp annotations.
At inference, both signals are composed via guided denoising, enabling annotation-free grasp synthesis for hands that participated \emph{only} in the morphology prior path.

Extensive experiments demonstrate that \modelname achieves state-of-the-art grasp synthesis across six robot hands spanning diverse morphologies, from 2-finger grippers to 21-link anthropomorphic hands.
We observe a clear positive \emph{scaling effect}: across all evaluated hands, per-hand grasp quality improves as the number of co-training embodiments increases, confirming that cross-embodiment diversity is a benefit rather than a compromise.
Furthermore, \modelname produces functional grasps for entirely new hands---given only a URDF and link meshes, with no object-grasp annotations---demonstrating strong annotation-free generalization.

\noindent In summary, our contributions are:
\begin{enumerate}
    \item We reformulate cross-embodiment grasp generation as a \emph{3D link-level assembly problem}, predicting per-link SE(3) poses in a kinematics-agnostic representation shared across hands of arbitrary topology---the enabling abstraction for a scalable, unified pipeline.
    \item We introduce a \emph{shared 3D shape-aware encoding} via a frozen pretrained VAE that maps both object surfaces and robot link meshes into a common geometric latent space, enabling the denoiser to reason jointly about object and link geometry without domain-specific encoders.
    \item We propose a \emph{classifier-free training paradigm} that learns a morphology prior from forward-kinematic self-exploration of valid hand configurations, enabling annotation-free transfer to novel robots given only a URDF and link point clouds.
    \item We achieve \emph{state-of-the-art} grasp synthesis across six morphologically diverse hands and reveal a positive scaling effect---individual hand performance improves with training embodiment diversity---while also demonstrating annotation-free generalization to unseen hands.
\end{enumerate}
